\lecture{15}{2026-11-05}{Provable Security}

\begin{itemize}
    \item Syntax: What your scheme/protocol should do. E.g., publick-key encryption, signature, etc.
    \item Goals: What security properties should the scheme satisfy?
    \item Instantiation: Describe specific algorithms matching the syntax (e.g. El-Gamal with elliptic curves).
    \item Hard Problem: e.g., DLOG, factoring, LWE.
    \item Proof: Prove that breaking the security definition is at least as hard as solving the hard problem.
\end{itemize}

\paragraph{Syntax.}
Public key encryption.
A PKE scheme is four sets \(\mathcal{PK}, \mathcal{SK}, \mathcal{M}, \mathcal{C}\) and three efficient algorithms 
\begin{align*} % add security parameter to keygen
    \mathsf{KeyGen} &: \text{randomness} \to \mathcal{PK} \times \mathcal{SK}, \\
    \mathsf{Enc} &: \mathcal{PK} \times \mathcal{M} \to \mathcal{C}, \\
    \mathsf{Dec} &: \mathcal{SK} \times \mathcal{C} \to \mathcal{M},
\end{align*}
or equivalently
\begin{align*}
    \mathsf{KeyGen}() & \to (\mathsf{pk}, \mathsf{sk}) \\
    \mathsf{Enc}(\mathsf{pk}, m) \to c \\
    \mathsf{Dec}(\mathsf{sk}, c) \to m.
\end{align*}

\paragraph{Goals.}
Correctness.
For all \((\mathsf{pk}, \mathsf{sk}) \in \mathcal{PK} \times \mathcal{SK}\) output by \(\mathsf{KeyGen}()\) and all \(m \in \mathcal{M}\), we have
\[
    \mathsf{Dec}(\mathsf{sk}, \mathsf{Enc}(\mathsf{pk}, m)) = m.
\]

Security.
One-wayness game:
\begin{itemize}
    \item \(m \xleftarrow{\$} \mathcal{M}\)
    \item \((\mathsf{pk}, \mathsf{sk}) \gets \mathsf{KeyGen}()\)
    \item \(c \gets \mathsf{Enc}(\mathsf{pk}, m)\)
    \item \(m' \gets \mathcal{A}(\mathsf{pk}, c)\)
    \item \(\mathcal{A}\) wins if \(m' = m\)
\end{itemize}

A PKE scheme is \emph{one-way secure} if for all efficient algorithms \(\mathcal{A}\), the probability that \(\mathcal{A}\) wins is negligible.

Efficient means probabilistic polynomial time (PPT) in the security parameter \(\lambda\). 
Negligible means that for all polynomials \(p(\cdot)\), for all sufficiently large \(\lambda\), we have \(\Pr[\mathcal{A} \text{ wins}] < 1/p(\lambda)\).

\paragraph{Instantiation.}
E.g. El-Gamal with elliptic curves.

\(\mathsf{KeyGen}()\):
\begin{itemize}
    \item \(x \xleftarrow{\$} \mathbb{Z}_q\)
    \item \(Q \gets xP\)
    \item return \((\mathsf{pk}, \mathsf{sk}) = (Q, x)\)
\end{itemize}

\(\mathsf{Enc}(\mathsf{pk}, m)\):
\begin{itemize}
    \item \(y \xleftarrow{\$} \mathbb{Z}_q\)
    \item \(S \gets y \mathsf{pk}\)
    \item \(c_1 \gets yP\)
    \item \(c_2 \gets m + S\)
    \item return \(c = (c_1, c_2)\)
\end{itemize}

\paragraph{Hard Problem.}
For El-Gamal with elliptic curves, the hard problem is the Computational Diffie-Hellman (CDH) problem: Given \(P, xP, yP\), compute \(xyP\).

\begin{definition}
    Problem \(A\) \emph{reduces to} problem \(B\) (denoted \(A \leq B\)) if there exists an efficient algorithm that solves \(A\) using an oracle for \(B\).
\end{definition}

\begin{proposition}
    \(\mathsf{CDH} \leq \mathsf{DLOG}\).
\end{proposition}
\begin{proof}
    Assume \(\mathcal{A}\) is an oracle that solves \(\mathsf{DLOG}\).
    Given an instance of \(\mathsf{CDH}\), i.e. \(P, xP, yP\), we can compute \(x = \mathcal{A}(P, xP)\) and \(y = \mathcal{A}(P, yP)\).
    Then we can compute \(xyP\) directly.
\end{proof}

\paragraph{Proof}
Almost always by contrapositive.
The real theorem is: If the hard problem is hard, then the scheme is secure.
The strategy to prove it is to assume the scheme is insecure, and show that this implies the hard problem is easy.

Concretely in our running example, we want to show that if El-Gamal is insecure, then \(\mathsf{CDH}\) is easy.

That is, for any efficient \(\mathcal{A}\) against OW-security of El-Gamal, we can construct an efficient \(\mathcal{S}\) that solves \(\mathsf{CDH}\).